\documentclass[11pt]{article}
\usepackage[a4paper,margin=1in]{geometry}
\usepackage{amsmath,amssymb,amsthm}

\theoremstyle{plain}
\newtheorem{theorem}{Theorem}[section]
\newtheorem{lemma}[theorem]{Lemma}
\newtheorem{proposition}[theorem]{Proposition}
\newtheorem{corollary}[theorem]{Corollary}
\theoremstyle{definition}
\newtheorem{definition}[theorem]{Definition}
\theoremstyle{remark}
\newtheorem{remark}[theorem]{Remark}
\newtheorem{example}[theorem]{Example}

% ------------------------------------------------------------------
% Notation (Shoenfield, Mathematical Logic, Ch.6, with code(u) for
% the encoding in place of corner-quotes).
% ------------------------------------------------------------------
\newcommand{\code}[1]{\mathit{code}(#1)}
\newcommand{\Pair}{\mathbf{Pair}}
\newcommand{\Vble}{\mathit{Vble}}
\newcommand{\Term}{\mathit{Term}}
\newcommand{\AFor}{\mathit{AFor}}
\newcommand{\For}{\mathit{For}}
\newcommand{\Fun}{\mathit{Fun}}
\newcommand{\Sub}{\mathit{Sub}}
\newcommand{\Subtl}{\mathit{Subtl}}
\newcommand{\Fr}{\mathit{Fr}}
\newcommand{\Defines}{\mathit{Defines}}
\newcommand{\Prf}{\mathit{Prf}}
\newcommand{\Prr}{\mathit{Pr}}
\newcommand{\Thmm}{\mathit{Thm}}
\newcommand{\Num}{\mathit{Num}}
\newcommand{\Con}{\mathit{Con}}
\newcommand{\lh}{\mathit{lh}}
\newcommand{\Len}{\mathit{Len}}
\newcommand{\SN}{\mathit{SN}}
\newcommand{\Eval}{\mathbf{Eval}}
\newcommand{\Berry}{\mathit{Berry}}
\newcommand{\bnt}{\mathbf{bin2nat}}
\newcommand{\bin}{\mathit{bin}}
\newcommand{\Bool}{\mathit{Bool}}
\newcommand{\cons}{\mathit{cons}}
\newcommand{\nilL}{\mathit{nil}}
\newcommand{\mul}{\mathbf{mul2}}
\newcommand{\add}{\mathbf{add}}

\title{A design note for a Berry-style proof of G\"odel's second incompleteness theorem for BRA\\
\large{(revised --- with binary-list compression of numerals)}}
\author{}
\date{}

\begin{document}
\maketitle

\begin{abstract}
This note develops a precise mathematical design for an alternative,
diagonal-free proof of G\"odel's second incompleteness theorem for the
basic recursive arithmetic~$T$ of Church~[1957] and Guard~[1963].
The argument is the user's proposal: take $\Fun_1$ itself as the
``\,naming language\,''; combine with Krivine's semantic uniformity
({\it Th\'eorie des ensembles}); and conclude via a Berry-style
self-application.  An earlier draft glossed over a critical
encoding-length issue: under Shoenfield's standard encoding, a numeral
$\mathbf{k}_n$ has $\Len(\code{\mathbf{k}_n}) = \Theta(n)$, defeating
the Berry inequality.  This revision introduces binary-list
compression of numerals via a fixed-size functional
$\bnt \in \Fun_1$, restoring $O(\log n)$ encoding for constants and
making the Berry contradiction bite at a concrete threshold.  We
estimate the implementation cost at approximately
$8\,000$ lines of Agda.  Throughout we follow Shoenfield's notation,
with $\code{u}$ for the encoding (in place of $\ulcorner u \urcorner$).
\end{abstract}

\tableofcontents

\section{Motivation: the two historical proofs of G\"odel II}
\label{sec:motivation}

There are two genuinely different proofs of the second incompleteness
theorem, and they were recognised as different by their authors at the
very start.  In his January 1931 letter to G\"odel --- written after
the K\"onigsberg meeting, conceding G\"odel's priority --- von Neumann
describes his own method:

\begin{quote}
``I believe I can reproduce your sequence of thoughts on the basis of
our communication and can therefore tell you that I used a somewhat
different method.  \emph{You prove $W \to A$, I show independently the
unprovability of $W$, though with a different kind of inference that
likewise copies the antinomies.}''
\end{quote}

Here $W$ is \emph{Widerspruchsfreiheit} (consistency, $\mathrm{Con}$)
and $A$ is the undecidable proposition (the G\"odel sentence $G$).  The
two methods are:

\begin{itemize}
\item \textbf{G\"odel's method ($W \to A$).}  Formalise G\"odel~I inside
the system to get $\vdash \mathrm{Con} \to G$; since $G$ is unprovable,
$\mathrm{Con}$ is unprovable.  This is the \emph{diagonal} route: build
the self-referential $G$, prove $\mathrm{Con}\to G$, combine with
G\"odel~I.

\item \textbf{von Neumann's method (direct, antinomy-copying).}  Show
\emph{independently} that $\mathrm{Con}$ is unprovable, by ``a different
kind of inference that likewise copies the antinomies'' --- i.e.\ by
\emph{directly} mimicking a paradox (the Richard/Berry/Liar family),
\emph{without} ever forming $G$ or proving $\mathrm{Con}\to G$.
\end{itemize}

This project realises both:

\begin{center}
\begin{tabular}{lll}
\hline
historical method & this project & engine \\
\hline
$W \to A$ (diagonal) & \texttt{BRA4.Thm.Thm14GodelII} (\emph{shipped}) & fixed point $+$ $\mathbf{sub}$ $+$ encoded mp \\
direct, antinomy-copying & the Berry route (\emph{this note}) & soundness $+$ Lindenbaum cut $+$ Berry counting \\
\hline
\end{tabular}
\end{center}

So the Berry route is not merely ``a second proof.''  Carried out
faithfully it would be the \textbf{von Neumann proof} standing beside
the \textbf{G\"odel proof} (Guard~1963, our diagonal development),
differing from it exactly as von Neumann told G\"odel they differ:
\emph{direct unprovability of $\mathrm{Con}$, copying an antinomy
(Berry), rather than the formalised $\mathrm{Con}\to G$.}  (von Neumann
never published his proof, so this is a faithful realisation of the
\emph{style} he describes, not a reconstruction of his actual argument.)

\section{Status, and the honest open problem}
\label{sec:status}

\textbf{This note's Theorem~\ref{thm:berry} (the original
``\,$T \not\vdash \Con_T$ via $\Sigma_0$-completeness\,'') is flawed as
written}, and the flaw is instructive.  It assumed a short Fun1
$\beta$ that \emph{Defines} $v(n_0)$ via the $\Eval$-based $\Sigma_0$
relation of Definition~\ref{def:defines}.  But to decide
$\nu_{n_0}(y)$, $\beta$ must run $\Eval$ on \emph{every} code of length
$< n_0$ --- including $\code{\beta}$ itself (since $\Len(\code\beta) <
n_0$).  As $\Eval(\code\beta, y) = \beta(y)$, this is an infinite
self-simulation regress: $\beta$ is not total, not a Fun1, and
``\,$\beta$ Defines $v(n_0)$\,'' is \emph{false} in $\mathbb{N}$ --- so
$\Sigma_0$-completeness cannot prove it and the contradiction never
fires.  This is the classical reason Berry's paradox is not an
outright contradiction.

\paragraph{What is correct, and the role of completeness.}  The fix is
not $\Eval$-naming ($\Sigma_0$) but \emph{provable-naming} ($\Pr_T$,
$\Sigma_1$): $v(n)$ must be defined through the provability predicate,
exactly as in Krivine.  The proof-search inherent in $\Pr_T$ is what
breaks the self-simulation regress.  Concretely the argument should be
\emph{von Neumann-style}: build, inside the syntactic
\emph{Tarski--Lindenbaum} universal model (Coquand--Smith constructive
completeness), a definability \emph{cut} that validates $T$ yet refutes
$\Con_T$; then \emph{soundness} (the easy direction) gives
$T \not\vdash \Con_T$.  In this framing the \emph{completeness theorem
proper is not load-bearing}: for an unprovability result the operative
half is soundness, and the ambient model is just the syntactic
Lindenbaum algebra (so ``definability $=$ provability'' is
definitional).  All the real content is the \textbf{cut construction
$+$ the Berry counting/length bound}, and the consumption of the
hypothesis $T \vdash \Con_T$ lives there.

\paragraph{What we do not yet have.}  We do \emph{not} have a concrete
construction of that cut for BRA.  The crux --- making the notion of
``definability/nameability'' precise enough that the Berry self-
reference produces a contradiction rather than dissolving --- is
exactly the difficulty flagged by van Dijk and Gietelink Oldenziel,
\emph{G\"odel's Incompleteness after Joyal} [2020], Remark~2.12 and its
footnote, in their arithmetic-universe account: \emph{``a more
conventional view isolates the vagueness in the notion of
`definability' or `nameability' as the culprit.''}  Their categorical proof (Joyal) sidesteps this by taking
the Cantor--Lawvere diagonal --- i.e.\ it is again the \emph{G\"odel}
($W\to A$) route, categorified in the predicative / formal-topology /
type-theoretic setting (the initial arithmetic universe $\mathcal U_0 =
\mathrm{Syn}(\mathcal T_{AU})$ is BRA's primitive-recursive world, the
categorical form of the Lindenbaum model).  So au.pdf \emph{confirms
the setting and names the difficulty, but does not provide the direct
construction.}

\paragraph{Update --- the cut is no longer vague.}  Reading Krivine's
proof \emph{as actually written} (\emph{Th\'eorie axiomatique des
ensembles}, pp.~126--127; see the companion note
\texttt{krivine-models-inside-models.tex}) replaces the unspecified
``Lindenbaum cut'' above with a concrete \emph{models-inside-models}
construction.  Krivine's uniformity is carried by: a model $M\models T$
which, under $T\vdash\Con_T$, contains an internal object $N'$ that $M$
believes is a model of~$T$, glued to $M$ by the soundness equivalence
``$N'\models F$ in $M$ iff $F\in T_M$.''  In BRA the outer model is
$\mathbf{th}$, the inner model is a $\mathbf{sub}$-relativisation of $\mathbf{th}$, and the
soundness link is $\mathbf{th}$'s validating-decoder invariant plus
$\Sigma_0$-completeness.  The hypothesis $T\vdash\Con_T$ is consumed at
exactly one place: the passage from $\mathbf{th}$ to its $\mathbf{sub}$-relativised
inner model.  This makes the ``cut'' a buildable target rather than an
open conceptual gap; only the Berry counting/length bound remains as
genuine (routine) work.

\paragraph{Consequence for the plan.}  Sections~\ref{sec:motivation}
onward describe the \emph{components} ($\Eval$, $\bnt$, $\Defines$,
$v$, the counting) faithfully; but the assembly into a theorem should
be read as the \emph{von Neumann / cut} target of {\S}\ref{sec:status},
with $\Pr_T$-naming replacing $\Eval$-naming, \emph{not} as the
$\Sigma_0$-only Theorem~\ref{thm:berry}.  The genuinely open step is
the cut construction; the genuinely hard but routine step is the Berry
counting/length bound (the recommended gate experiment).

\section{Setting and notation}

\subsection{The theory $T$ and the grammars $\Fun_1, \Fun_2$}

We write $T$ for the basic recursive arithmetic of Church~[1957], in
the formulation of Guard~[1963].  Its non-logical symbols are a
constant $\mathbf{O}$, a unary symbol $\mathbf{s}$, the Church
combinators $\mathbf{o}, \mathbf{u}, \mathbf{v}, \mathbf{C},
\mathbf{R}$, and the binary identity~$=$.  Variables are $\mathbf{x}_0,
\mathbf{x}_1, \ldots$ with $\SN(\mathbf{x}_i) = 2i$.

The two functional grammars are
$$
\mathbf{f} ::= \mathbf{s} \mid \mathbf{o} \mid \mathbf{u} \mid \mathbf{C}(\mathbf{g}, \mathbf{f_1}, \mathbf{f_2}) \quad (\Fun_1),
\qquad
\mathbf{g} ::= \mathbf{v} \mid \mathbf{R}(\mathbf{f}, \mathbf{g_1}, \mathbf{g_2}) \quad (\Fun_2).
$$
Every $\mathbf{f} \in \Fun_1$ denotes a total recursive function
$\mathbb{N} \to \mathbb{N}$.  We write $\vdash \mathbf{A}$ for
``\,$\mathbf{A} \in \Thmm_T$\,''.

\subsection{Numerals, encoding, and length}

The numeral is $\mathbf{k}_n := \mathbf{s}^n(\mathbf{O})$.  The
encoding $\code{u}$ is defined as in {\tt goedelII-summary.tex} via
Pair-based sequences.  We write $\Len(a)$ for the symbol count of the
encoding tree of~$a$:
$$
\Len(\code{\mathbf{O}}) = 1, \quad
\Len(\code{\mathbf{s}}) = 1, \quad
\Len(\code{\mathbf{x}_i}) = O(1), \quad
\Len(\code{\mathbf{u}(\mathbf{v_1}, \ldots, \mathbf{v_n})}) = 1 + \sum_i \Len(\code{\mathbf{v_i}}),
$$
and analogously for the other formation rules.  In particular,
\begin{equation}
\Len(\code{\mathbf{k}_n}) = n + 1, \quad \text{(numerals have linear encoding length).}
\label{eq:numeral-linear}
\end{equation}
This linear growth is the obstacle the binary-list compression below
is designed to overcome.

\section{Internal evaluation of $\Fun_1$}

\begin{lemma}\label{lem:eval}
There is a binary functional $\Eval \in \Fun_2$ such that for every
$\mathbf{f} \in \Fun_1$ and every closed numeric term~$\mathbf{t}$,
$$
\vdash \Eval(\Num(\code{\mathbf{f}}), \mathbf{t}) = \mathbf{f}(\mathbf{t}).
$$
\end{lemma}

\noindent
$\Eval$ is constructed by course-of-values recursion on the four
$\Fun_1$ tags ($\mathbf{s}, \mathbf{o}, \mathbf{u}, \mathbf{C}$) via
the standard Church combinators, in the same style as the verifier
$\mathbf{thmT}$ of \verb|BRA4/ThmT.agda|, but with the verifier-cascade
reduced from seventeen Hilbert-axiom tags plus three rule tags to just
four functional-grammar tags (plus a single dispatch tag for the
binary sub-grammar $\Fun_2$).  Estimated cost in the Agda
implementation:
\begin{itemize}
\item $\Eval$ definition + cov-spec dispatch + base/step closures: $\sim 800$ LoC;
\item per-tag closure proofs ($\Eval$ at $\mathbf{s}, \mathbf{o}, \mathbf{u}, \mathbf{C}, \mathbf{v}, \mathbf{R}$): $\sim 1200$ LoC.
\end{itemize}

\section{Binary-list compression of numerals}
\label{sec:bin}

\subsection{The encoding}

We encode a meta-natural~$n$ as a finite list of bits, least-significant
first.  Bits are $\mathbf{O}$ (for $0$) and $\mathbf{s}\,\mathbf{O}$
(for $1$).  Lists are encoded via $\Pair$ with $\mathbf{O}$ as
``\,$\nilL$\,''.  Define meta-recursively
\begin{equation}
\bin(0) := \mathbf{O}, \qquad
\bin(n) := \cons\bigl(b_n,\, \bin\bigl(\lfloor n/2 \rfloor\bigr)\bigr)
\quad (n > 0),
\label{eq:bin-meta}
\end{equation}
where $b_n = \mathbf{O}$ if $n$ is even and $b_n = \mathbf{s}\,\mathbf{O}$
otherwise, and $\cons(b, \ell) := \Pair(b, \ell)$.

\begin{lemma}\label{lem:bin-length}
For every $n \geq 1$, $\Len(\code{\bin(n)}) = O(\log n)$.
\end{lemma}

\begin{proof}
Each bit contributes a constant amount; the list has
$\lceil \log_2(n+1) \rceil$ bits.
\end{proof}

\subsection{Decoding back to a numeral}

\begin{definition}\label{def:bin2nat}
Define $\bnt \in \Fun_1$ by primitive recursion on the list shape:
\begin{align*}
\bnt(\mathbf{O}) &= \mathbf{O}, \\
\bnt(\cons(\mathbf{b}, \ell)) &= \add(\mathbf{b}, \mul(\bnt(\ell))),
\end{align*}
where $\mul, \add \in \Fun_1$ are the doubling and successor-of-double
functionals, each of fixed code length.
\end{definition}

\begin{remark}
Strictly, $\bnt$ recurses on a $\Pair$-encoded list, which fits the
$\Fun_2$-via-$\mathbf{R}$ pattern of Church.  We elide the standard
gymnastics of fitting list-recursion into $\Fun_1$ syntax; what
matters is that $\bnt$ has a fixed Agda code length, independent of
its argument.
\end{remark}

\begin{lemma}\label{lem:bin-correct}
For every $n \in \mathbb{N}$, $\vdash \bnt(\bin(n)) = \mathbf{k}_n$.
\end{lemma}

\begin{proof}[Proof sketch]
Meta-induction on the number of bits of $n$.  Base case $n = 0$:
$\vdash \bnt(\mathbf{O}) = \mathbf{O}$ by definition.  Step: assume
$\vdash \bnt(\bin(\lfloor n/2 \rfloor)) = \mathbf{k}_{\lfloor n/2
\rfloor}$.  Then by the closure on $\cons$ and the closures of $\mul,
\add$ on closed numerals,
$$
\vdash \bnt(\bin(n)) = \add(b_n, \mul(\mathbf{k}_{\lfloor n/2 \rfloor})) = \mathbf{k}_n.
$$
The number of derivation steps is $O(\log n)$, hence so is the
encoding length of the proof.
\end{proof}

\begin{corollary}\label{cor:bin-len}
For every $n$, the closed term $\bnt(\bin(n))$ has
$\Len(\code{\bnt(\bin(n))}) \leq C + d \log_2(n + 1)$ for fixed
meta-constants $C, d$ (depending on $|\code{\bnt}|$ and the cons-cell
overhead, but not on~$n$).
\end{corollary}

\noindent
By contrast, the standard numeral $\mathbf{k}_n$ has
$\Len(\code{\mathbf{k}_n}) = n + 1$ by
\eqref{eq:numeral-linear}.  The compression factor is exponential.

\subsection{Substituting the compressed form into formulas}

Wherever a meta-natural~$n$ enters a formula or term of~$T$ in the
sequel, we use the compressed form $\bnt(\bin(n))$ in place of
$\mathbf{k}_n$.  By Lemma~\ref{lem:bin-correct}, all such uses are
provably equal to the corresponding $\mathbf{k}_n$, so the formula
they appear in is preserved under $\vdash$-equivalence; but the
\emph{encoding length} of the formula is reduced from linear in~$n$
to logarithmic in~$n$.

\section{The $\Defines$ predicate}

\begin{definition}\label{def:defines}
At the meta level:
$$
\Defines(a, b) \;\leftrightarrow\;
\Eval(a, \mathbf{k}_b) \neq 0 \;\wedge\; \forall c_{c < b}.\,\Eval(a, \mathbf{k}_c) = 0,
$$
restricted to $a$ being the encoding of a $\Fun_1$ functional.
\end{definition}

\noindent
Internally, $\Defines$ is realised by a $\Fun_2$ test
$\mathbf{TestDef}$ such that
$$
\vdash \mathbf{TestDef}(\mathbf{a}, \mathbf{b}) = \mathbf{O} \;\Longleftrightarrow\; \neg \Defines(\mathrm{val}\,\mathbf{a}, \mathrm{val}\,\mathbf{b})
\quad\text{in $\mathbb{N}$,}
$$
where $\mathrm{val}$ is meta-evaluation.  $\mathbf{TestDef}$ is built
from $\Eval$ plus a small bounded-recursion helper for the $\forall c
< b$ clause, all of fixed Agda code length (estimated $\sim 400$ LoC).

For closed inputs $\mathbf{a} = \bnt(\bin(\code{\mathbf{f}}))$ and
$\mathbf{b} = \bnt(\bin(y))$, the equation
$\mathbf{TestDef}(\mathbf{a}, \mathbf{b}) = \mathbf{O}$ is a $\Sigma_0$
sentence which is true iff $\mathbf{f}$ does not define~$y$, and which
$T$ proves by direct evaluation through the $\bnt, \Eval$ closures.

\section{The Krivine-Berry function $v$}

\begin{definition}\label{def:vfn}
For $n \in \mathbb{N}$,
$$
v(n) \;:=\; \mu y \in \mathbb{N}.\ \bigl(\,\forall a \in \mathbb{N}.\,\Len(a) < n \;\Rightarrow\; \neg \Defines(a, y)\,\bigr).
$$
\end{definition}

\begin{lemma}\label{lem:vtotal}
$v$ is a total primitive-recursive function $\mathbb{N} \to \mathbb{N}$.
\end{lemma}

\begin{proof}
The set $\{a : \Len(a) < n\}$ is finite with bound $K_n \leq U(n)$ for
some explicit primitive-recursive bound $U$.  Each $a$ defines at most
one $y$, so the set of $y$'s defined by some such $a$ has cardinality
$\leq K_n$; its complement is non-empty.  $v(n)$ is the least element
of this complement, computed by bounded search in primitive recursion.
\end{proof}

\section{Internalising $v$ via a compact predicate}

\subsection{The predicate $\nu_n$}

For each meta-natural $n$ we define
$$
\nu_n(\mathbf{x}_0) \;\equiv\; \bigl(\, \mathbf{TestAllDef}(\bnt(\bin(n)),\, \mathbf{x}_0) = \mathbf{O} \,\bigr),
$$
where $\mathbf{TestAllDef} \in \Fun_2$ is the bounded-universal
helper that returns $\mathbf{O}$ iff no code of $\Len < n$ defines its
second argument.  Concretely:
$$
\mathbf{TestAllDef}(\mathbf{n}, \mathbf{y}) := \bigl[\,\forall a_{a \leq U(\mathbf{n})}.\,(\Len(a) < \mathbf{n} \;\supset\; \mathbf{TestDef}(a, \mathbf{y}) = \mathbf{O})\,\bigr]\text{-witness Fun},
$$
implemented as a primitive recursion on $U(\mathbf{n})$ with bounded
existential check.  $\mathbf{TestAllDef}$ has fixed Agda code length.

\begin{lemma}\label{lem:nu-length}
$\Len(\code{\nu_n(\mathbf{x}_0)}) \leq C' + d' \log_2(n + 1)$ for
meta-constants $C', d'$ depending on $|\code{\mathbf{TestAllDef}}|$ and
$|\code{\bnt}|$, but not on~$n$.
\end{lemma}

\begin{proof}
The formula is $\mathbf{TestAllDef}(\bnt(\bin(n)), \mathbf{x}_0) =
\mathbf{O}$.  Its encoding tree has fixed top structure
$\Pair(\mathbf{k}_{\mathrm{tag\_eq}}, \Pair(\cdot, \mathbf{O}))$, with
the only $n$-dependent part being the closed term
$\bnt(\bin(n))$, which by Corollary~\ref{cor:bin-len} has length $C +
d \log_2(n + 1)$.  Hence
$\Len(\code{\nu_n}) = C' + d' \log_2(n+1)$ with $C' = O(C) + |\code{\mathbf{TestAllDef}}|
+ |\code{\mathbf{x}_0}| + \mathrm{const}$.
\end{proof}

\subsection{The $\Sigma_0$-completeness chain}
\label{sec:sigma0}

\begin{proposition}[$\Sigma_0$-completeness for compressed-numeral Eval-equations]\label{prop:sound}
For every $\mathbf{f} \in \Fun_1$ and $b \in \mathbb{N}$, if
$\mathbf{f}(b) = c$ in the standard model, then
$$
\vdash \Eval(\bnt(\bin(\code{\mathbf{f}})), \bnt(\bin(b))) = \bnt(\bin(c)).
$$
The encoded derivation has $\Len$ equal to $O(\Len(\code{\mathbf{f}})
+ \log b + \log c)$.
\end{proposition}

\begin{proof}[Proof sketch]
Compose three derivations:
\begin{enumerate}
\item $\vdash \bnt(\bin(\code{\mathbf{f}})) = \mathbf{k}_{\code{\mathbf{f}}}$ and $\vdash \bnt(\bin(b)) = \mathbf{k}_b$ by Lemma~\ref{lem:bin-correct}.
\item $\vdash \Eval(\mathbf{k}_{\code{\mathbf{f}}}, \mathbf{k}_b) = \mathbf{k}_c$ by Lemma~\ref{lem:eval} (the verifier closure on the un-compressed numeral).
\item $\vdash \mathbf{k}_c = \bnt(\bin(c))$ by Lemma~\ref{lem:bin-correct} applied symmetrically.
\end{enumerate}
The first and third have $\Len$ $O(\log b)$ and $O(\log c)$
respectively (by following the binary expansions).  The middle
derivation has $\Len$ proportional to the number of computation steps
of $\mathbf{f}$ on input $b$, which itself is bounded by a function
primitive-recursive in $\Len(\code{\mathbf{f}})$ and $b$.

(For our applications, $\mathbf{f}$ is always a fixed small functional
$\mathbf{TestDef}, \mathbf{TestAllDef}$, etc., so this middle term is
$O(\log b)$ uniformly.)
\end{proof}

\begin{corollary}\label{cor:nu-sound}
If $\vdash \nu_n(\bnt(\bin(y)))$, then $v(n) \leq y$, i.e.\ $y$ is not
defined by any code of length $< n$.  Conversely, for the specific
value $y = v(n)$, $\vdash \nu_n(\bnt(\bin(v(n))))$ is derivable in
$O(\Len(\code{\nu_n}) + \log v(n)) = O(\log n)$ encoded
derivation steps.
\end{corollary}

\section{The Berry contradiction}

\subsection{The witness formula $\beta$}

Define
$$
\beta(\mathbf{x}_0) \;:=\; \nu_{n_0}(\mathbf{x}_0) \;\wedge\; \mathbf{TestNoneLess}(\bnt(\bin(n_0)),\, \mathbf{x}_0) = \mathbf{O},
$$
where $\mathbf{TestNoneLess}(\mathbf{n}, \mathbf{x}_0)$ is the
$\Fun_2$ predicate ``\,every $\mathbf{x}_1 < \mathbf{x}_0$ has
$\nu_{\mathbf{n}}(\mathbf{x}_1)$ false\,'', realised as a bounded
primitive recursion on $\mathbf{x}_0$.

\begin{lemma}\label{lem:beta-len}
$\Len(\code{\beta(\mathbf{x}_0)}) \leq C'' + d'' \log_2(n_0 + 1)$.
\end{lemma}

\begin{proof}
$\beta$ is the conjunction of two atomic equations, each of the form
``\,$F(\bnt(\bin(n_0)), \mathbf{x}_0) = \mathbf{O}$\,'' with $F$ a
fixed functional.  Apply Lemma~\ref{lem:nu-length} and analogously for
$\mathbf{TestNoneLess}$.
\end{proof}

\begin{lemma}\label{lem:beta-defines}
In the standard model, $\beta(y)$ holds iff $y = v(n_0)$.  Hence
$$
\Defines(\code{\beta}, v(n_0))
\quad\text{holds in $\mathbb{N}$,}
$$
where $\code{\beta}$ is interpreted as the unary $\Fun_1$
characteristic-function of $\beta$.
\end{lemma}

\subsection{The contradiction}

\begin{theorem}\label{thm:berry}
$T \not\vdash \Con_T$.
\end{theorem}

\begin{proof}[Proof sketch]
Suppose for contradiction $\vdash \Con_T$.

Pick~$n_0$ such that $C'' + d'' \log_2(n_0 + 1) < n_0$ (possible by
Lemma~\ref{lem:beta-len} for any sufficiently large $n_0$; concretely
$n_0 \sim 10^3$ suffices for typical $C'', d''$).

By Lemma~\ref{lem:beta-defines} and the proof of
Lemma~\ref{lem:vtotal}, $v(n_0)$ is in fact defined (in the
$\Defines$-sense of Definition~\ref{def:defines}) by~$\code{\beta}$.
By Proposition~\ref{prop:sound} (which uses $\vdash \Con_T$ only to
exclude the bottom case --- otherwise the $\Sigma_0$-completeness step
is unconditional), $\vdash \mathbf{TestDef}(\bnt(\bin(\code{\beta})),
\bnt(\bin(v(n_0)))) \neq \mathbf{O}$ and the bounded-zero-below
condition holds.

But $\Len(\code{\beta}) < n_0$, and by the definition of $v(n_0)$, no
code of length $< n_0$ defines $v(n_0)$.  Contradiction.
\end{proof}

\begin{remark}\label{rem:role-of-con}
The assumption $\vdash \Con_T$ is used precisely once: to ensure that
the formal derivation of the $\Sigma_0$ sentence ``\,$\code{\beta}$
defines $v(n_0)$\,'' actually goes through (rather than $T$
deriving its negation as well, which would be the contradiction we
want).  This is the analogue of Krivine's uniformity-across-models, in
single-model form: the standard model $\mathbb{N}$ satisfies the
$\Sigma_0$ sentence, and consistency lifts truth to provability.  No
H\"ilbert--Bernays--L\"ob conditions, no diagonal lemma.
\end{remark}

\section{Honest accounting: what we still need}

\subsection{What is avoided}

The Berry route avoids the following components of the diagonal proof
in \verb|BRA4/Thm/Thm14GodelII.agda|:
\begin{itemize}
\item the encoded substitution functors $\mathbf{sub}, \mathbf{sub_2},
\mathbf{sub_3}$ and their per-shape closure equations
($\sim 10\,000$ LoC eliminated);
\item the predicate-Leibniz lemma $\mathbf{substF\_cong}$
(BRA3.Substitutivity);
\item the Carneiro-lifted hypothetical modus-ponens helper toolkit
$\mathbf{imp\_encoded\_mp}$, $\mathbf{impEqTrans}$, etc.;
\item the inlined re-derivation of G\"odel~I for the sub-form of $G$;
\item the diagonal lemma itself.
\end{itemize}

\subsection{What is still needed}

\begin{itemize}
\item $\Eval \in \Fun_2$: a verifier for $\Fun_1$, structurally
analogous to $\mathbf{thmT}$ but with only four functional-tag
closures.  This is a non-trivial development ($\sim 2000$ LoC), albeit
smaller than the full $\mathbf{thmT}$ cascade.

\item $\bnt \in \Fun_1$ and the soundness chain through it: the
compact-numeral functor of Section~\ref{sec:bin}.  About $\sim 800$
LoC: definition, correctness lemma (\ref{lem:bin-correct}), and the
$\Sigma_0$-completeness composition of Proposition~\ref{prop:sound}.

\item The decidable test functionals $\mathbf{TestDef},
\mathbf{TestAllDef}, \mathbf{TestNoneLess}$ and their closure
equations: $\sim 1500$ LoC.

\item The meta-level $v$ and totality by counting: $\sim 300$ LoC.

\item The witness formula $\beta$, its length bound, and the Berry
contradiction Theorem~\ref{thm:berry}: $\sim 1500$ LoC, including the
$O(\log n)$ derivation chains through $\bnt$.

\item Concrete computation of constants $C', C'', d', d''$ and a
specific witness $n_0$: $\sim 500$ LoC.  This is the part that requires
careful counting; the proof in principle works for any sufficiently
large~$n_0$, but to be useful as an effective bound the constants
must be made explicit.
\end{itemize}

\paragraph{Revised total estimate.}  About $\sim 8\,000$ LoC, comparable
to but slightly smaller than the diagonal-based proof's
$\sim 12\,000$ LoC.  The earlier figure of $\sim 5\,500$ LoC was over-optimistic
because it underweighted the binary-compression infrastructure (about
$\sim 800$ LoC) and the $\Sigma_0$-completeness chain through $\bnt$
(about $\sim 1\,000$ LoC).

\section{Comparison with the diagonal proof}

\begin{center}
\begin{tabular}{p{6cm} c c}
\hline
& Diagonal & Berry/Krivine \\
\hline
Verifier $\mathbf{thmT}$ for derivations & full cascade & --- \\
$\Eval$ for $\Fun_1$ & --- & required \\
Encoded substitution ($\mathbf{sub}$, $\mathbf{sub_2}$) & required & --- \\
Diagonal lemma & required & --- \\
$\mathbf{substF\_cong}$ (predicate-Leibniz) & required & --- \\
Carneiro lift $\mathbf{imp\_encoded\_mp}$ & required & --- \\
Inlined G\"odel I (sub-form) & required & --- \\
$\Sigma_0$-completeness for $\Eval$-equations & --- & required \\
Binary-list numeral compression $\bnt$ & --- & required \\
Bounded-universal helpers $\mathbf{TestAllDef}$ etc. & --- & required \\
Meta-level Berry function $v$ & --- & required \\
Encoding-length budget $\Len(\code{\beta}) < n_0$ & --- & required, concrete \\
\hline
Estimated LoC & $\sim 12\,000$ & $\sim 8\,000$ \\
\hline
\end{tabular}
\end{center}

\section{The role of $\Fun_1$ as ``naming language''}

The user's idea of using $\Fun_1$ itself as the naming language has
three concrete advantages over the formulations of Boolos and Chaitin
discussed in Berry.pdf:

\begin{enumerate}
\item \textbf{Built-in totality.}  Every $\mathbf{f} \in \Fun_1$
denotes a \emph{total} recursive function, by Church's construction.
``\,$\mathbf{f}$ defines $y$\,'' is an unconditional statement
involving $\Eval$, with no totality side condition.

\item \textbf{Concrete recursive codes.}  $\code{\mathbf{f}}$ is a
structurally-recursive encoding tree.  $\Len$ is recursive and
well-behaved, and the binary-list compression of
Section~\ref{sec:bin} works uniformly through it.

\item \textbf{Symmetric ``\,defines\,''.}  Boolos's ``\,$\varphi$
names $n$\,'' bundles a $\Sigma_1$-witness with a $\Pi_1$-uniqueness
clause, requiring $\Sigma_1$-completeness in its full form.  The
$\Fun_1$ version $\Defines(\code{\mathbf{f}}, y)$ replaces this by the
conjunction of a positive value with a bounded
zero-below-condition --- both bounded, both recursive, both
$\Sigma_0$-completeness suffices.
\end{enumerate}

\section{Remark on $\mathbf{sub_2}$}

A side benefit of the Berry route: the development entirely avoids the
two-variable simultaneous substitution functor $\mathbf{sub_2}$ which
was used as a convenience in the diagonal-based proof
(\verb|BRA4/SbT2.agda| and its dependents, $\sim 10\,000$ LoC).  As
noted in \verb|BRA4/goedelII-summary.tex| {\S}5, $\mathbf{sub_2}$
admits a factorisation into three single-variable substitutions, but
this factorisation was deferred and never implemented in the diagonal
route.  In the Berry route the issue does not arise at all, since the
construction uses only $\Eval$ (single-variable in its proof index)
and the per-tag $\Fun_1$ functional helpers.

\section{Summary}

The Krivine-style Berry argument, instantiated with $\Fun_1$ as the
naming language and binary-list compression of numerals
(Section~\ref{sec:bin}), gives a diagonal-free proof of G\"odel's
second incompleteness theorem for~$T$, requiring:
\begin{itemize}
\item a small verifier $\Eval \in \Fun_2$ for $\Fun_1$;
\item the compact-numeral functor $\bnt \in \Fun_1$ with
$\Len(\code{\bnt(\bin(n))}) = O(\log n)$;
\item the decidable $\Defines$ predicate and bounded-universal helpers;
\item $\Sigma_0$-completeness for $\Eval$-equations through
$\bnt$; and
\item the compressed witness formula $\beta$ of length
$O(\log n_0)$ for the contradiction.
\end{itemize}
The argument bypasses the diagonal lemma, the
H\"ilbert--Bernays--L\"ob conditions, the substitutivity lemma
$\mathbf{substF\_cong}$, the two-variable substitution machinery
$\mathbf{sub_2}$, and the encoded modus-ponens Carneiro lift, all of
which were used in the diagonal-based proof already shipped in
\verb|BRA4/Thm/Thm14GodelII.agda|.  The estimated implementation cost
is $\sim 8\,000$ LoC.

\paragraph{Honest caveat.}  The Berry route is \emph{not} a fundamentally
cheaper alternative; it is an \emph{architecturally different}
alternative.  Its principal scientific value is to provide a second,
independent proof of G\"odel~II for~$T$ with no shared infrastructure
beyond Church~1957's $\Fun_1, \Fun_2$ grammars and the basic Hilbert
calculus.  The savings are real but more modest than a naive count
suggests, because the binary-list compression machinery
($\bnt$, $\bin$, soundness chain) costs $\sim 1\,800$ LoC of its own.

\end{document}
